\subsection{Scope of the Results}
Three classes of statement are supported. The analytical class:
Propositions~\ref{prop:exact} and~\ref{prop:multi} are exact under their stated
assumptions, and the figures accompanying Proposition~\ref{prop:approx} are maxima
observed on the stated grid, properties of the idealised mixture rather than of any
fleet. The simulated class: everything in Section~\ref{sec:design} is generated
under~\eqref{eq:sim} with assumed parameters. The measured class: only the capacitor
figures of Section~\ref{sec:evidence}. A manifest accompanying the code gives every
reported value's script, result file and provenance tag; no tag denotes a shipboard
measurement, because there is none.

\subsection{Limitations}
\subsubsection{No Shipboard Quantity}
No failure rate, temperature coefficient, variance
component, coverage or spares count for shipboard power electronics is reported,
because none could be estimated. The values $\sigma=\mpSigma$ and $\tau,\omega$
used in Sections~\ref{sec:multi} and~\ref{sec:design} illustrate the arithmetic
and must not be quoted as fleet characteristics.

\subsubsection{No Empirical Calibration}
No predictive interval has been checked against
observed shipboard failures, so the calibration of the predictive distribution
remains open and is the first thing a campaign should test.

\subsubsection{No Identified Dependence Structure}
The division of latent spread between
the ship-wide and interaction terms is not determined by any evidence here. Since
it changes a multi-reference plan materially under either criterion
(Section~\ref{sec:multi}) and is difficult to estimate
(Section~\ref{sec:design}), this limits any provisioning recommendation,
including the one developed here.

\subsubsection{Idealised Theory}
Propositions~\ref{prop:exact} and~\ref{prop:approx}
concern the mixed-Poisson case~\eqref{eq:mixture}, not the full model with its
cell dispersion, crossed effects and parameter uncertainty; the common-cause and
cascading mechanisms of Section~\ref{sec:mechanisms} are outside
Assumption~\ref{as:mix} altogether and act on the upper tail, where a
provisioning quantile lies. A plan computed under Assumption~\ref{as:mix} will
understate the stock required if such events occur at non-negligible rates.

\subsubsection{No Remaining-Life Estimate}
The model gives a failure count over a bounded
horizon; extending it to remaining life requires unit-level time-to-event records
and a survival or recurrent-event formulation~\cite{meeker2022} within the same
hierarchy. Finally, the federated protocol is not the cheaper option for
cell-aggregated counts and offers no confidentiality guarantee
(Section~\ref{sec:fed}).

\subsection{Requirements for Empirical Validation}
\label{sec:needed}
Six records are necessary. The first is an inventory of the installed population
by ship, zone, reference, manufacturer and installation date, maintained through
refits. The second is a set of event records in which every replacement is dated
and classified as functional failure or preventive by a rule fixed in advance,
with removals for refit, obsolescence and cannibalisation recorded separately;
where that classification is ambiguous the ambiguity must be quantified and the
analysis repeated under both readings. The third is exposure in powered
unit-hours per population and mission. The fourth is a set of covariates logged
over complete missions, so that the previous mission's summaries exist at the
moment of decision, since a covariate available only contemporaneously cannot
enter this decision however predictive it is. The fifth is repeated missions and
several references per ship, without which the two terms of~\eqref{eq:structure}
cannot be separated (Section~\ref{sec:design}). The sixth is a record of
common-cause events, together with teardown of a sample of removed units, so
that conditional independence and the limiting failure mode are tested rather
than assumed.

With such records the analysis developed here can be carried out on ships, and its
calibration, variance components and dependence structure become measurable. In
their absence those quantities cannot be estimated, and none is reported.
